\section{Design Patterns \& Resilience}
{
    \metroset{titleformat frame=smallcaps}

% --- Slide 1 (words) ---
\begin{frame}{Designing for Failure: Why Resilience Matters}
\begin{itemize}
    \item In a distributed system, failures are normal: networks blip, services restart, requests get duplicated.
    \item A naive application crashes or corrupts data when this happens.
    \item \textcolor{HUAblue}{\textbf{The idea:}} cloud-native design patterns turn these failures into handled, recoverable events.
\end{itemize}
\end{frame}

% --- Slide 2 (words) ---
\begin{frame}{Decoupling \& Recovery}
\begin{itemize}
    \item \textbf{Claim Check:} large product images go to MinIO; only a lightweight reference travels through RabbitMQ.
    \item \textbf{Retry (exponential backoff):} transient failures when talking to RabbitMQ and MongoDB are retried automatically.
    \item \textbf{Circuit Breaker:} if a dependency is down, the breaker opens and fails fast, preventing cascading failures.
\end{itemize}
\end{frame}

% --- Slide 3 (words) ---
\begin{frame}{Safety \& Scale}
\begin{itemize}
    \item \textbf{Idempotency:} an idempotency key prevents duplicate products on a retried or double-submitted request.
    \item \textbf{Stateless services:} the backend keeps no local state -- everything lives in MongoDB/MinIO -- so it scales horizontally and survives restarts.
    \item \textcolor{HUAblue}{\textbf{Together:}} these patterns make the system fault-tolerant and ready to scale.
\end{itemize}
\end{frame}

}


